\chapter{Firmware}

The main file \texttt{Pinacoteca.ino} instantiates modules and defines pins. The logic is organized in two distinct phases, providing a predictable and easily verifiable execution cycle.

\section{Setup}
During \texttt{setup()} the following are initialized:
\begin{itemize}
  \item serial for debug and diagnostics;
  \item servo and turnstile with a known initial position;
  \item traffic light, thermostat, lighting, and humidity;
  \item LCD panel with a startup message.
\end{itemize}
Initializing the servo before other modules reduces the risk of unintended movement during pin configuration.

\section{Loop}
In the main loop, each module updates its state in a coherent sequence:
\begin{enumerate}
  \item update turnstile and people count;
  \item update traffic light based on the limit;
  \item adjust climate (heating or cooling);
  \item adjust lighting;
  \item control humidity;
  \item refresh the display with current data.
\end{enumerate}
The order is intentional: human inputs are handled first, then status indicators are updated, and finally the LCD interface is refreshed. This way the display reflects the latest decision of the cycle.

\section{Error handling}
Modules that read sensors validate measurement correctness. If input is invalid, outputs are set to a safe state and the update function returns \texttt{false}. This behavior aligns with the project fail-safe philosophy.

\section{Target configuration}
Targets are passed to module constructors:
\begin{itemize}
  \item \texttt{Thermostat}: desired temperature in Celsius.
  \item \texttt{LightingControl}: target illuminance in lux.
  \item \texttt{HumidifierControl}: target humidity in percent.
  \item \texttt{Turnstile}: maximum number of people.
\end{itemize}
Centralized configuration in the main file allows quick tuning without modifying individual modules.
